\documentclass[11pt, a4paper]{article}
\usepackage[a4paper, margin=5.0cm]{geometry}
% \usepackage{fontspec}
\usepackage[spanish]{babel}
% \babelprovide[import, onchar=ids fonts]{spanish}
% \babelfont{rm}{Noto Sans}
\usepackage{amsmath}
\usepackage{amssymb}
\usepackage{enumitem}

\pagestyle{empty}

\begin{document}

\noindent
\textbf{Ejercicio:} Un calorímetro ideal y adiabático contiene, inicialmente, una masa de aceite de oliva a $80^\circ\text{C}$. Se introduce en el calorímetro un bloque de $100\text{ g}$ de hielo a $0^\circ\text{C}$. Luego de un rato, se encuentra que el hielo se ha fundido por completo, y tanto el agua como el aceite se encuentran a $40^\circ\text{C}$. Sabiendo que el calor específico del aceite de oliva es $c_\text{aceite} = 0.4\textrm{\,cal/(g $^\circ$C)}$, determinar la masa de aceite de oliva contenida en el calorímetro.

\end{document}